% =========================================================
% comandos.tex
% Colección de comandos personalizados organizados por tipo
% =========================================================
%
% Este fichero define atajos para:
% 1) Citas y notas al pie
% 2) Cabeceras
% 3) Enlaces y pies de figura
% 4) Etiquetas visuales de respuesta
% 5) Ajustes globales
% 6) Índices y glosario
%
% Recomendación:
% Mantén este archivo cargado en el preámbulo con:
% % =========================================================
% comandos.tex
% Colección de comandos personalizados organizados por tipo
% =========================================================
%
% Este fichero define atajos para:
% 1) Citas y notas al pie
% 2) Cabeceras
% 3) Enlaces y pies de figura
% 4) Etiquetas visuales de respuesta
% 5) Ajustes globales
% 6) Índices y glosario
%
% Recomendación:
% Mantén este archivo cargado en el preámbulo con:
% % =========================================================
% comandos.tex
% Colección de comandos personalizados organizados por tipo
% =========================================================
%
% Este fichero define atajos para:
% 1) Citas y notas al pie
% 2) Cabeceras
% 3) Enlaces y pies de figura
% 4) Etiquetas visuales de respuesta
% 5) Ajustes globales
% 6) Índices y glosario
%
% Recomendación:
% Mantén este archivo cargado en el preámbulo con:
% % =========================================================
% comandos.tex
% Colección de comandos personalizados organizados por tipo
% =========================================================
%
% Este fichero define atajos para:
% 1) Citas y notas al pie
% 2) Cabeceras
% 3) Enlaces y pies de figura
% 4) Etiquetas visuales de respuesta
% 5) Ajustes globales
% 6) Índices y glosario
%
% Recomendación:
% Mantén este archivo cargado en el preámbulo con:
% \input{comandos}
%
% ---------------------------------------------------------
% 1. CITAS Y NOTAS AL PIE
% ---------------------------------------------------------
%
% Estos comandos ayudan a insertar citas bibliográficas o web
% en notas al pie de forma uniforme.
%
% Requisitos habituales:
% - biblatex
% - \fullcite disponible
%
% Convención:
% - #1, #2, #3... representan los parámetros del comando
%

% \footcitepage{texto citado}{clave_biblatex}{pagina}
% Sirve para mostrar un fragmento entre comillas y añadir en nota al pie
% la referencia completa indicando la página concreta.
%
% Ejemplo:
% \footcitepage{La seguridad es esencial}{nist2024}{15}
%
% Resultado:
% “La seguridad es esencial”¹
% y en la nota al pie: referencia completa, p. 15
\newcommand{\footcitepage}[3]{%
  ``#1''\footnote{%
    \fullcite{#2}, p.~#3%
  }%
}

% \footbibliographypage{clave_biblatex}{pagina}
% Sirve para insertar solo una nota al pie con la referencia completa
% y la página, sin mostrar texto citado en el cuerpo del documento.
%
% Ejemplo:
% \footbibliographypage{nist2024}{15}
\newcommand{\footbibliographypage}[2]{%
  \footnote{%
    \fullcite{#1}, p.~#2%
  }%
}

% \footwebfragment{clave_biblatex}{fragmento}
% Sirve para citar una fuente web en nota al pie e indicar qué fragmento
% concreto de esa fuente se está usando.
%
% Ejemplo:
% \footwebfragment{owasp2024}{Se recomienda validar todas las entradas}
\newcommand{\footwebfragment}[2]{%
  \footnote{%
    \fullcite{#1}. Fragmento: #2%
  }%
}

% \footweb{clave_biblatex}
% Sirve para insertar una nota al pie con la referencia completa de una web
% o cualquier otra fuente, sin añadir fragmento ni página.
%
% Ejemplo:
% \footweb{owasp2024}
\newcommand{\footweb}[1]{%
  \footnote{%
    \fullcite{#1}%
  }%
}

% \footwebcite{texto citado}{clave_biblatex}{fragmento}
% Sirve para mostrar un texto entre comillas en el cuerpo del documento
% y añadir en nota al pie la fuente web más el fragmento relevante.
%
% Ejemplo:
% \footwebcite{La autenticación multifactor reduce riesgos}{cisa2024}{MFA helps prevent unauthorized access}
\newcommand{\footwebcite}[3]{%
  ``#1''\footnote{%
    \fullcite{#2}. Fragmento: #3%
  }%
}

% \footcitedefinition{termino_o_definicion}{clave_biblatex}{pagina}
% Sirve para escribir un término o definición y añadir la fuente en nota
% al pie con la página correspondiente.
%
% Ejemplo:
% \footcitedefinition{Confidencialidad}{iso27001}{8}
\newcommand{\footcitedefinition}[3]{%
  #1\footnote{%
    \fullcite{#2}, p.~#3%
  }%
}

% ---------------------------------------------------------
% 2. CABECERAS
% ---------------------------------------------------------
%
% Estos comandos están pensados para usar con fancyhdr.
% Permiten personalizar el contenido de la cabecera.
%
% Posiciones típicas en fancyhdr:
% - L = izquierda
% - C = centro
% - R = derecha
% - E = página par
% - O = página impar
%
% Ejemplos de posiciones:
% - LE, RO, RH, LH, etc.
%

% \nombreCapituloCabecera{posicion}{texto}
% Sirve para poner un texto en cursiva en una posición concreta
% de la cabecera.
%
% Ejemplo:
% \nombreCapituloCabecera{RH}{Capítulo 1}
\newcommand{\nombreCapituloCabecera}[2]{%
  \fancyhead[#1]{\textit{#2}}%
}

% \logocabeceraCEU{posicion}{ruta_imagen}
% Sirve para colocar un logotipo o imagen en la cabecera.
%
% Ejemplo:
% \logocabeceraCEU{LH}{Imagenes/logo-ceu.png}
%
% Nota:
% La altura de la imagen está fijada a 1 cm.
\newcommand{\logocabeceraCEU}[2]{%
  \fancyhead[#1]{\includegraphics[height=1cm]{#2}}%
}

% ---------------------------------------------------------
% 3. ENLACES, FIGURAS Y ELEMENTOS VISUALES
% ---------------------------------------------------------

% \enlace{url}{texto_visible}
% Sirve para crear un hipervínculo con color personalizado y subrayado.
%
% Ejemplo:
% \enlace{https://www.example.com}{Sitio oficial}
%
% Requisitos habituales:
% - hyperref
% - ulem o un paquete que defina \uline
% - color/bluelink definido en el documento
\newcommand{\enlace}[2]{%
  \href{#1}{\color{bluelink}\uline{#2}}%
}

% \captionazul{texto_caption}
% Sirve para crear una caption de figura o tabla con texto en color azul.
%
% Ejemplo:
% \begin{figure}
%   ...
%   \captionazul{Arquitectura general del sistema}
% \end{figure}
%
% Nota:
% También usa el mismo texto para la entrada corta del índice de figuras/tablas.
\newcommand{\captionazul}[1]{%
  \caption[#1]{\textcolor{blueColor}{#1}}%
}

% ---------------------------------------------------------
% 4. ETIQUETAS VISUALES DE RESPUESTA
% ---------------------------------------------------------
%
% Estos comandos imprimen respuestas cortas en color.
% Son útiles en tablas comparativas o matrices de decisión.
%

% \permits
% Muestra “Sí” en verde.
%
% Ejemplo:
% Acceso de administrador: \permits
\newcommand{\permits}{\textcolor{green!60!black}{Sí}}

% \denies
% Muestra “No” en rojo.
%
% Ejemplo:
% Acceso de invitado: \denies
\newcommand{\denies}{\textcolor{red!70!black}{No}}

% ---------------------------------------------------------
% 5. AJUSTES GLOBALES
% ---------------------------------------------------------
%
% Estas líneas modifican comportamientos por defecto de paquetes.
%

% Elimina la línea horizontal de la cabecera definida por fancyhdr.
\renewcommand{\headrulewidth}{0pt}

% Evita que aparezcan números de página junto a las entradas del glosario.
% Suele usarse cuando no interesa mostrar en qué páginas aparece cada término.
\renewcommand*{\glossaryentrynumbers}[1]{}

% ---------------------------------------------------------
% 6. ÍNDICES Y GLOSARIO
% ---------------------------------------------------------
%
% Estos comandos generan índices y glosario con ajustes de espaciado.
%

% \crearIndiceTOC
% Genera el índice general (table of contents).
%
% Ejemplo:
% \crearIndiceTOC
\newcommand{\crearIndiceTOC}{%
    \vspace*{-2.5cm}%
    \tableofcontents%
}

% \crearIndiceTablas
% Genera el índice de tablas.
%
% Ejemplo:
% \crearIndiceTablas
\newcommand{\crearIndiceTablas}{%
    \vspace*{-2.5cm}%
    \listoftables%
}

% \crearIndiceFiguras
% Genera el índice de figuras.
%
% Ejemplo:
% \crearIndiceFiguras
\newcommand{\crearIndiceFiguras}{%
    \vspace*{-2.5cm}%
    \listoffigures%
}

% \crearglosario
% Inserta el glosario de términos, ajusta la cabecera y añade todas
% las entradas del glosario aunque no hayan sido usadas en el texto.
%
% Ejemplo:
% \crearglosario
%
% Nota:
% Este comando asume que existe el fichero:
% Glosario/glosario.tex
\newcommand{\crearglosario}{%
    \vspace*{-1cm}%
    \nombreCapituloCabecera{RH}{Glosario de términos}%

    \glsaddall%
    \vspace*{-2cm}%
    \input{Glosario/glosario}%
}
%
% ---------------------------------------------------------
% 1. CITAS Y NOTAS AL PIE
% ---------------------------------------------------------
%
% Estos comandos ayudan a insertar citas bibliográficas o web
% en notas al pie de forma uniforme.
%
% Requisitos habituales:
% - biblatex
% - \fullcite disponible
%
% Convención:
% - #1, #2, #3... representan los parámetros del comando
%

% \footcitepage{texto citado}{clave_biblatex}{pagina}
% Sirve para mostrar un fragmento entre comillas y añadir en nota al pie
% la referencia completa indicando la página concreta.
%
% Ejemplo:
% \footcitepage{La seguridad es esencial}{nist2024}{15}
%
% Resultado:
% “La seguridad es esencial”¹
% y en la nota al pie: referencia completa, p. 15
\newcommand{\footcitepage}[3]{%
  ``#1''\footnote{%
    \fullcite{#2}, p.~#3%
  }%
}

% \footbibliographypage{clave_biblatex}{pagina}
% Sirve para insertar solo una nota al pie con la referencia completa
% y la página, sin mostrar texto citado en el cuerpo del documento.
%
% Ejemplo:
% \footbibliographypage{nist2024}{15}
\newcommand{\footbibliographypage}[2]{%
  \footnote{%
    \fullcite{#1}, p.~#2%
  }%
}

% \footwebfragment{clave_biblatex}{fragmento}
% Sirve para citar una fuente web en nota al pie e indicar qué fragmento
% concreto de esa fuente se está usando.
%
% Ejemplo:
% \footwebfragment{owasp2024}{Se recomienda validar todas las entradas}
\newcommand{\footwebfragment}[2]{%
  \footnote{%
    \fullcite{#1}. Fragmento: #2%
  }%
}

% \footweb{clave_biblatex}
% Sirve para insertar una nota al pie con la referencia completa de una web
% o cualquier otra fuente, sin añadir fragmento ni página.
%
% Ejemplo:
% \footweb{owasp2024}
\newcommand{\footweb}[1]{%
  \footnote{%
    \fullcite{#1}%
  }%
}

% \footwebcite{texto citado}{clave_biblatex}{fragmento}
% Sirve para mostrar un texto entre comillas en el cuerpo del documento
% y añadir en nota al pie la fuente web más el fragmento relevante.
%
% Ejemplo:
% \footwebcite{La autenticación multifactor reduce riesgos}{cisa2024}{MFA helps prevent unauthorized access}
\newcommand{\footwebcite}[3]{%
  ``#1''\footnote{%
    \fullcite{#2}. Fragmento: #3%
  }%
}

% \footcitedefinition{termino_o_definicion}{clave_biblatex}{pagina}
% Sirve para escribir un término o definición y añadir la fuente en nota
% al pie con la página correspondiente.
%
% Ejemplo:
% \footcitedefinition{Confidencialidad}{iso27001}{8}
\newcommand{\footcitedefinition}[3]{%
  #1\footnote{%
    \fullcite{#2}, p.~#3%
  }%
}

% ---------------------------------------------------------
% 2. CABECERAS
% ---------------------------------------------------------
%
% Estos comandos están pensados para usar con fancyhdr.
% Permiten personalizar el contenido de la cabecera.
%
% Posiciones típicas en fancyhdr:
% - L = izquierda
% - C = centro
% - R = derecha
% - E = página par
% - O = página impar
%
% Ejemplos de posiciones:
% - LE, RO, RH, LH, etc.
%

% \nombreCapituloCabecera{posicion}{texto}
% Sirve para poner un texto en cursiva en una posición concreta
% de la cabecera.
%
% Ejemplo:
% \nombreCapituloCabecera{RH}{Capítulo 1}
\newcommand{\nombreCapituloCabecera}[2]{%
  \fancyhead[#1]{\textit{#2}}%
}

% \logocabeceraCEU{posicion}{ruta_imagen}
% Sirve para colocar un logotipo o imagen en la cabecera.
%
% Ejemplo:
% \logocabeceraCEU{LH}{Imagenes/logo-ceu.png}
%
% Nota:
% La altura de la imagen está fijada a 1 cm.
\newcommand{\logocabeceraCEU}[2]{%
  \fancyhead[#1]{\includegraphics[height=1cm]{#2}}%
}

% ---------------------------------------------------------
% 3. ENLACES, FIGURAS Y ELEMENTOS VISUALES
% ---------------------------------------------------------

% \enlace{url}{texto_visible}
% Sirve para crear un hipervínculo con color personalizado y subrayado.
%
% Ejemplo:
% \enlace{https://www.example.com}{Sitio oficial}
%
% Requisitos habituales:
% - hyperref
% - ulem o un paquete que defina \uline
% - color/bluelink definido en el documento
\newcommand{\enlace}[2]{%
  \href{#1}{\color{bluelink}\uline{#2}}%
}

% \captionazul{texto_caption}
% Sirve para crear una caption de figura o tabla con texto en color azul.
%
% Ejemplo:
% \begin{figure}
%   ...
%   \captionazul{Arquitectura general del sistema}
% \end{figure}
%
% Nota:
% También usa el mismo texto para la entrada corta del índice de figuras/tablas.
\newcommand{\captionazul}[1]{%
  \caption[#1]{\textcolor{blueColor}{#1}}%
}

% ---------------------------------------------------------
% 4. ETIQUETAS VISUALES DE RESPUESTA
% ---------------------------------------------------------
%
% Estos comandos imprimen respuestas cortas en color.
% Son útiles en tablas comparativas o matrices de decisión.
%

% \permits
% Muestra “Sí” en verde.
%
% Ejemplo:
% Acceso de administrador: \permits
\newcommand{\permits}{\textcolor{green!60!black}{Sí}}

% \denies
% Muestra “No” en rojo.
%
% Ejemplo:
% Acceso de invitado: \denies
\newcommand{\denies}{\textcolor{red!70!black}{No}}

% ---------------------------------------------------------
% 5. AJUSTES GLOBALES
% ---------------------------------------------------------
%
% Estas líneas modifican comportamientos por defecto de paquetes.
%

% Elimina la línea horizontal de la cabecera definida por fancyhdr.
\renewcommand{\headrulewidth}{0pt}

% Evita que aparezcan números de página junto a las entradas del glosario.
% Suele usarse cuando no interesa mostrar en qué páginas aparece cada término.
\renewcommand*{\glossaryentrynumbers}[1]{}

% ---------------------------------------------------------
% 6. ÍNDICES Y GLOSARIO
% ---------------------------------------------------------
%
% Estos comandos generan índices y glosario con ajustes de espaciado.
%

% \crearIndiceTOC
% Genera el índice general (table of contents).
%
% Ejemplo:
% \crearIndiceTOC
\newcommand{\crearIndiceTOC}{%
    \vspace*{-2.5cm}%
    \tableofcontents%
}

% \crearIndiceTablas
% Genera el índice de tablas.
%
% Ejemplo:
% \crearIndiceTablas
\newcommand{\crearIndiceTablas}{%
    \vspace*{-2.5cm}%
    \listoftables%
}

% \crearIndiceFiguras
% Genera el índice de figuras.
%
% Ejemplo:
% \crearIndiceFiguras
\newcommand{\crearIndiceFiguras}{%
    \vspace*{-2.5cm}%
    \listoffigures%
}

% \crearglosario
% Inserta el glosario de términos, ajusta la cabecera y añade todas
% las entradas del glosario aunque no hayan sido usadas en el texto.
%
% Ejemplo:
% \crearglosario
%
% Nota:
% Este comando asume que existe el fichero:
% Glosario/glosario.tex
\newcommand{\crearglosario}{%
    \vspace*{-1cm}%
    \nombreCapituloCabecera{RH}{Glosario de términos}%

    \glsaddall%
    \vspace*{-2cm}%
    \input{Glosario/glosario}%
}
%
% ---------------------------------------------------------
% 1. CITAS Y NOTAS AL PIE
% ---------------------------------------------------------
%
% Estos comandos ayudan a insertar citas bibliográficas o web
% en notas al pie de forma uniforme.
%
% Requisitos habituales:
% - biblatex
% - \fullcite disponible
%
% Convención:
% - #1, #2, #3... representan los parámetros del comando
%

% \footcitepage{texto citado}{clave_biblatex}{pagina}
% Sirve para mostrar un fragmento entre comillas y añadir en nota al pie
% la referencia completa indicando la página concreta.
%
% Ejemplo:
% \footcitepage{La seguridad es esencial}{nist2024}{15}
%
% Resultado:
% “La seguridad es esencial”¹
% y en la nota al pie: referencia completa, p. 15
\newcommand{\footcitepage}[3]{%
  ``#1''\footnote{%
    \fullcite{#2}, p.~#3%
  }%
}

% \footbibliographypage{clave_biblatex}{pagina}
% Sirve para insertar solo una nota al pie con la referencia completa
% y la página, sin mostrar texto citado en el cuerpo del documento.
%
% Ejemplo:
% \footbibliographypage{nist2024}{15}
\newcommand{\footbibliographypage}[2]{%
  \footnote{%
    \fullcite{#1}, p.~#2%
  }%
}

% \footwebfragment{clave_biblatex}{fragmento}
% Sirve para citar una fuente web en nota al pie e indicar qué fragmento
% concreto de esa fuente se está usando.
%
% Ejemplo:
% \footwebfragment{owasp2024}{Se recomienda validar todas las entradas}
\newcommand{\footwebfragment}[2]{%
  \footnote{%
    \fullcite{#1}. Fragmento: #2%
  }%
}

% \footweb{clave_biblatex}
% Sirve para insertar una nota al pie con la referencia completa de una web
% o cualquier otra fuente, sin añadir fragmento ni página.
%
% Ejemplo:
% \footweb{owasp2024}
\newcommand{\footweb}[1]{%
  \footnote{%
    \fullcite{#1}%
  }%
}

% \footwebcite{texto citado}{clave_biblatex}{fragmento}
% Sirve para mostrar un texto entre comillas en el cuerpo del documento
% y añadir en nota al pie la fuente web más el fragmento relevante.
%
% Ejemplo:
% \footwebcite{La autenticación multifactor reduce riesgos}{cisa2024}{MFA helps prevent unauthorized access}
\newcommand{\footwebcite}[3]{%
  ``#1''\footnote{%
    \fullcite{#2}. Fragmento: #3%
  }%
}

% \footcitedefinition{termino_o_definicion}{clave_biblatex}{pagina}
% Sirve para escribir un término o definición y añadir la fuente en nota
% al pie con la página correspondiente.
%
% Ejemplo:
% \footcitedefinition{Confidencialidad}{iso27001}{8}
\newcommand{\footcitedefinition}[3]{%
  #1\footnote{%
    \fullcite{#2}, p.~#3%
  }%
}

% ---------------------------------------------------------
% 2. CABECERAS
% ---------------------------------------------------------
%
% Estos comandos están pensados para usar con fancyhdr.
% Permiten personalizar el contenido de la cabecera.
%
% Posiciones típicas en fancyhdr:
% - L = izquierda
% - C = centro
% - R = derecha
% - E = página par
% - O = página impar
%
% Ejemplos de posiciones:
% - LE, RO, RH, LH, etc.
%

% \nombreCapituloCabecera{posicion}{texto}
% Sirve para poner un texto en cursiva en una posición concreta
% de la cabecera.
%
% Ejemplo:
% \nombreCapituloCabecera{RH}{Capítulo 1}
\newcommand{\nombreCapituloCabecera}[2]{%
  \fancyhead[#1]{\textit{#2}}%
}

% \logocabeceraCEU{posicion}{ruta_imagen}
% Sirve para colocar un logotipo o imagen en la cabecera.
%
% Ejemplo:
% \logocabeceraCEU{LH}{Imagenes/logo-ceu.png}
%
% Nota:
% La altura de la imagen está fijada a 1 cm.
\newcommand{\logocabeceraCEU}[2]{%
  \fancyhead[#1]{\includegraphics[height=1cm]{#2}}%
}

% ---------------------------------------------------------
% 3. ENLACES, FIGURAS Y ELEMENTOS VISUALES
% ---------------------------------------------------------

% \enlace{url}{texto_visible}
% Sirve para crear un hipervínculo con color personalizado y subrayado.
%
% Ejemplo:
% \enlace{https://www.example.com}{Sitio oficial}
%
% Requisitos habituales:
% - hyperref
% - ulem o un paquete que defina \uline
% - color/bluelink definido en el documento
\newcommand{\enlace}[2]{%
  \href{#1}{\color{bluelink}\uline{#2}}%
}

% \captionazul{texto_caption}
% Sirve para crear una caption de figura o tabla con texto en color azul.
%
% Ejemplo:
% \begin{figure}
%   ...
%   \captionazul{Arquitectura general del sistema}
% \end{figure}
%
% Nota:
% También usa el mismo texto para la entrada corta del índice de figuras/tablas.
\newcommand{\captionazul}[1]{%
  \caption[#1]{\textcolor{blueColor}{#1}}%
}

% ---------------------------------------------------------
% 4. ETIQUETAS VISUALES DE RESPUESTA
% ---------------------------------------------------------
%
% Estos comandos imprimen respuestas cortas en color.
% Son útiles en tablas comparativas o matrices de decisión.
%

% \permits
% Muestra “Sí” en verde.
%
% Ejemplo:
% Acceso de administrador: \permits
\newcommand{\permits}{\textcolor{green!60!black}{Sí}}

% \denies
% Muestra “No” en rojo.
%
% Ejemplo:
% Acceso de invitado: \denies
\newcommand{\denies}{\textcolor{red!70!black}{No}}

% ---------------------------------------------------------
% 5. AJUSTES GLOBALES
% ---------------------------------------------------------
%
% Estas líneas modifican comportamientos por defecto de paquetes.
%

% Elimina la línea horizontal de la cabecera definida por fancyhdr.
\renewcommand{\headrulewidth}{0pt}

% Evita que aparezcan números de página junto a las entradas del glosario.
% Suele usarse cuando no interesa mostrar en qué páginas aparece cada término.
\renewcommand*{\glossaryentrynumbers}[1]{}

% ---------------------------------------------------------
% 6. ÍNDICES Y GLOSARIO
% ---------------------------------------------------------
%
% Estos comandos generan índices y glosario con ajustes de espaciado.
%

% \crearIndiceTOC
% Genera el índice general (table of contents).
%
% Ejemplo:
% \crearIndiceTOC
\newcommand{\crearIndiceTOC}{%
    \vspace*{-2.5cm}%
    \tableofcontents%
}

% \crearIndiceTablas
% Genera el índice de tablas.
%
% Ejemplo:
% \crearIndiceTablas
\newcommand{\crearIndiceTablas}{%
    \vspace*{-2.5cm}%
    \listoftables%
}

% \crearIndiceFiguras
% Genera el índice de figuras.
%
% Ejemplo:
% \crearIndiceFiguras
\newcommand{\crearIndiceFiguras}{%
    \vspace*{-2.5cm}%
    \listoffigures%
}

% \crearglosario
% Inserta el glosario de términos, ajusta la cabecera y añade todas
% las entradas del glosario aunque no hayan sido usadas en el texto.
%
% Ejemplo:
% \crearglosario
%
% Nota:
% Este comando asume que existe el fichero:
% Glosario/glosario.tex
\newcommand{\crearglosario}{%
    \vspace*{-1cm}%
    \nombreCapituloCabecera{RH}{Glosario de términos}%

    \glsaddall%
    \vspace*{-2cm}%
    \input{Glosario/glosario}%
}
%
% ---------------------------------------------------------
% 1. CITAS Y NOTAS AL PIE
% ---------------------------------------------------------
%
% Estos comandos ayudan a insertar citas bibliográficas o web
% en notas al pie de forma uniforme.
%
% Requisitos habituales:
% - biblatex
% - \fullcite disponible
%
% Convención:
% - #1, #2, #3... representan los parámetros del comando
%

% \footcitepage{texto citado}{clave_biblatex}{pagina}
% Sirve para mostrar un fragmento entre comillas y añadir en nota al pie
% la referencia completa indicando la página concreta.
%
% Ejemplo:
% \footcitepage{La seguridad es esencial}{nist2024}{15}
%
% Resultado:
% “La seguridad es esencial”¹
% y en la nota al pie: referencia completa, p. 15
\newcommand{\footcitepage}[3]{%
  ``#1''\footnote{%
    \fullcite{#2}, p.~#3%
  }%
}

% \footbibliographypage{clave_biblatex}{pagina}
% Sirve para insertar solo una nota al pie con la referencia completa
% y la página, sin mostrar texto citado en el cuerpo del documento.
%
% Ejemplo:
% \footbibliographypage{nist2024}{15}
\newcommand{\footbibliographypage}[2]{%
  \footnote{%
    \fullcite{#1}, p.~#2%
  }%
}

% \footwebfragment{clave_biblatex}{fragmento}
% Sirve para citar una fuente web en nota al pie e indicar qué fragmento
% concreto de esa fuente se está usando.
%
% Ejemplo:
% \footwebfragment{owasp2024}{Se recomienda validar todas las entradas}
\newcommand{\footwebfragment}[2]{%
  \footnote{%
    \fullcite{#1}. Fragmento: #2%
  }%
}

% \footweb{clave_biblatex}
% Sirve para insertar una nota al pie con la referencia completa de una web
% o cualquier otra fuente, sin añadir fragmento ni página.
%
% Ejemplo:
% \footweb{owasp2024}
\newcommand{\footweb}[1]{%
  \footnote{%
    \fullcite{#1}%
  }%
}

% \footwebcite{texto citado}{clave_biblatex}{fragmento}
% Sirve para mostrar un texto entre comillas en el cuerpo del documento
% y añadir en nota al pie la fuente web más el fragmento relevante.
%
% Ejemplo:
% \footwebcite{La autenticación multifactor reduce riesgos}{cisa2024}{MFA helps prevent unauthorized access}
\newcommand{\footwebcite}[3]{%
  ``#1''\footnote{%
    \fullcite{#2}. Fragmento: #3%
  }%
}

% \footcitedefinition{termino_o_definicion}{clave_biblatex}{pagina}
% Sirve para escribir un término o definición y añadir la fuente en nota
% al pie con la página correspondiente.
%
% Ejemplo:
% \footcitedefinition{Confidencialidad}{iso27001}{8}
\newcommand{\footcitedefinition}[3]{%
  #1\footnote{%
    \fullcite{#2}, p.~#3%
  }%
}

% ---------------------------------------------------------
% 2. CABECERAS
% ---------------------------------------------------------
%
% Estos comandos están pensados para usar con fancyhdr.
% Permiten personalizar el contenido de la cabecera.
%
% Posiciones típicas en fancyhdr:
% - L = izquierda
% - C = centro
% - R = derecha
% - E = página par
% - O = página impar
%
% Ejemplos de posiciones:
% - LE, RO, RH, LH, etc.
%

% \nombreCapituloCabecera{posicion}{texto}
% Sirve para poner un texto en cursiva en una posición concreta
% de la cabecera.
%
% Ejemplo:
% \nombreCapituloCabecera{RH}{Capítulo 1}
\newcommand{\nombreCapituloCabecera}[2]{%
  \fancyhead[#1]{\textit{#2}}%
}

% \logocabeceraCEU{posicion}{ruta_imagen}
% Sirve para colocar un logotipo o imagen en la cabecera.
%
% Ejemplo:
% \logocabeceraCEU{LH}{Imagenes/logo-ceu.png}
%
% Nota:
% La altura de la imagen está fijada a 1 cm.
\newcommand{\logocabeceraCEU}[2]{%
  \fancyhead[#1]{\includegraphics[height=1cm]{#2}}%
}

% ---------------------------------------------------------
% 3. ENLACES, FIGURAS Y ELEMENTOS VISUALES
% ---------------------------------------------------------

% \enlace{url}{texto_visible}
% Sirve para crear un hipervínculo con color personalizado y subrayado.
%
% Ejemplo:
% \enlace{https://www.example.com}{Sitio oficial}
%
% Requisitos habituales:
% - hyperref
% - ulem o un paquete que defina \uline
% - color/bluelink definido en el documento
\newcommand{\enlace}[2]{%
  \href{#1}{\color{bluelink}\uline{#2}}%
}

% \captionazul{texto_caption}
% Sirve para crear una caption de figura o tabla con texto en color azul.
%
% Ejemplo:
% \begin{figure}
%   ...
%   \captionazul{Arquitectura general del sistema}
% \end{figure}
%
% Nota:
% También usa el mismo texto para la entrada corta del índice de figuras/tablas.
\newcommand{\captionazul}[1]{%
  \caption[#1]{\textcolor{blueColor}{#1}}%
}

% ---------------------------------------------------------
% 4. ETIQUETAS VISUALES DE RESPUESTA
% ---------------------------------------------------------
%
% Estos comandos imprimen respuestas cortas en color.
% Son útiles en tablas comparativas o matrices de decisión.
%

% \permits
% Muestra “Sí” en verde.
%
% Ejemplo:
% Acceso de administrador: \permits
\newcommand{\permits}{\textcolor{green!60!black}{Sí}}

% \denies
% Muestra “No” en rojo.
%
% Ejemplo:
% Acceso de invitado: \denies
\newcommand{\denies}{\textcolor{red!70!black}{No}}

% ---------------------------------------------------------
% 5. AJUSTES GLOBALES
% ---------------------------------------------------------
%
% Estas líneas modifican comportamientos por defecto de paquetes.
%

% Elimina la línea horizontal de la cabecera definida por fancyhdr.
\renewcommand{\headrulewidth}{0pt}

% Evita que aparezcan números de página junto a las entradas del glosario.
% Suele usarse cuando no interesa mostrar en qué páginas aparece cada término.
\renewcommand*{\glossaryentrynumbers}[1]{}

% ---------------------------------------------------------
% 6. ÍNDICES Y GLOSARIO
% ---------------------------------------------------------
%
% Estos comandos generan índices y glosario con ajustes de espaciado.
%

% \crearIndiceTOC
% Genera el índice general (table of contents).
%
% Ejemplo:
% \crearIndiceTOC
\newcommand{\crearIndiceTOC}{%
    \vspace*{-2.5cm}%
    \tableofcontents%
}

% \crearIndiceTablas
% Genera el índice de tablas.
%
% Ejemplo:
% \crearIndiceTablas
\newcommand{\crearIndiceTablas}{%
    \vspace*{-2.5cm}%
    \listoftables%
}

% \crearIndiceFiguras
% Genera el índice de figuras.
%
% Ejemplo:
% \crearIndiceFiguras
\newcommand{\crearIndiceFiguras}{%
    \vspace*{-2.5cm}%
    \listoffigures%
}

% \crearglosario
% Inserta el glosario de términos, ajusta la cabecera y añade todas
% las entradas del glosario aunque no hayan sido usadas en el texto.
%
% Ejemplo:
% \crearglosario
%
% Nota:
% Este comando asume que existe el fichero:
% Glosario/glosario.tex
\newcommand{\crearglosario}{%
    \vspace*{-1cm}%
    \nombreCapituloCabecera{RH}{Glosario de términos}%

    \glsaddall%
    \vspace*{-2cm}%
    \input{Glosario/glosario}%
}